\chapter{对偶}\label{V}
\origpage{160}{154}
本篇汇集与\term{Poincaré 对偶}{Poincaré duality}{dualité de Poincaré}有关的几个定理和相容性。第 \ref{V:1} 节讨论一维局部双对偶定理（\cite[I 5.1]{sga5}），并作若干对偶的计算。第 \ref{V:2} 节介绍 M.\ Artin 告诉我的一个十分简洁的曲线上 Poincaré 对偶的证明。第 \ref{V:3} 节证明一个相容性，将给出曲线上 Poincaré 对偶的配对的两种定义联系起来：一种通过杯积，另一种通过 Jacobi 簇的自身对偶。最后，第 \ref{V:4} 节证明其标题中所指的相容性。

本篇的概形均为分离诺特概形，$n$ 为在所考虑的每个概形上可逆的整数。

\section{一维局部双对偶}\label{V:1}
\origpage{161}{155}
\subsection{}\label{V:1-1}
设 $S$ 为纯维数 $1$ 的正则概形。我们要证明集中于次数 $0$ 的 $\dZ/n$ 复形为\term{对偶化复形}{dualizing complex}{complexe dualisant}，即：对 $\sK\in\ob\eD_c^b(S,\dZ/n)$（下标 $c$ 表示可构造），若置 $D\sK=\rHom(\sK,\dZ/n)$，则 $D\sK$ 仍有可构造的有界上同调，且从 $\sK$ 到 $DD\sK$ 的典范态射 $\alpha$ 为同构。

对 $\sK\in\eD^-$ 及任意 $\sL$，有
\[
\hom(\sK\lotimes\sL,\dZ/n)=\hom(\sL,\rHom(\sK,\dZ/n)).
\]
在假定 $D\sK\in\eD^-$ 后，便可定义 $\sK\to DD\sK$：若 $\beta:D\sK\lotimes\sK\to\dZ/n$ 为典范配对，则此态射由配对
$\sK\lotimes D\sK=D\sK\lotimes\sK\xrightarrow\beta\dZ/n$ 定义。

\subsection{}\label{V:1-2}
设 $f:X\to S$ 为分离有限型态射，置 $\sK_X=\eR f^!\dZ/n$。对 $\sK\in\ob\eD^-(X,\dZ/n)$，置 $D\sK=\rHom(\sK,\sK_X)$。由 $\eR f_!$ 与 $\eR f^!$ 的伴随关系，
\[
\eR f_*D\sK=\eR f_*\rHom(\sK,\eR f^!\dZ/n)
=\rHom(\eR f_!\sK,\dZ/n)=D\eR f_!\sK.
\]
$\eR f_!\sK$ 与 $\eR f_*D\sK$ 的对应配对为复合
\[
\eR f_!\sK\lotimes\eR f_*D\sK\longrightarrow\eR f_!(\sK\lotimes D\sK)
\longrightarrow\eR f_!\eR f^!\dZ/n\longrightarrow\dZ/n.
\]
这一描述的对称性表明，当 $f$ 固有时，下图交换：
\begin{equation*}\tag{1.2.1}\label{V:eq:1-2-1}
\xymatrix{\eR f_*\sK\ar[r]\ar[d]&DD\eR f_*\sK\ar[d]^-\wr\\
\eR f_*DD\sK\ar[r]^-\sim&D\eR f_*D\sK.}
\end{equation*}

若 $f$ 为一个闭点的包含，已知 $\eR f^!\dZ/n=\dZ/n(-1)[-2]$，即除扭曲与平移外就是 $\dZ/n$。在域的谱上，该复形为对偶化复形（这是 $\dZ/n$-模的 Pontrjagin 对偶）。由 \eqref{V:eq:1-2-1}，对 $\sK=\eR f_!\sL$ 形式的对象有 $\sK\iso DD\sK$；从而当 $\underline\h^\bullet(\sK)$ 的支撑有限时，亦有此结论。

\begin{theorem_}\label{V:1-3}
设 $j:U\hookrightarrow S$ 为稠密开集，$\sF$ 为 $U$ 上局部常可构造的 $\dZ/n$-模层。则 $Dj_*\sF=j_*D\sF$，即
\[
\Hom(j_*\sF,\dZ/n)=j_*\Hom(\sF,\dZ/n),\qquad
\Ext^i(j_*\sF,\dZ/n)=0\quad(i>0).
\]
\end{theorem_}
在 $U$ 上，由于 $\sF$ 局部常且 $\dZ/n$ 是内射 $\dZ/n$-模，$i>0$ 时 $\Ext^i(j_*\sF,\dZ/n)=0$；$i=0$ 时则为 $\sF$ 的对偶 $\sF^\vee$。可直接验证 $\Hom(j_*\sF,\dZ/n)=j_*\sF^\vee$，故只须在 $S\setminus U$ 的各点验证 $i>0$ 时 $\Ext^i$ 为零。问题在这些点处是局部的，因而可假定 $S$ 为严格局部\term{性状}{trait}{trait}，$U$ 仅由其泛点 $\eta$ 组成。设 $I=\gal(\bar\eta/\eta)$。层 $\sF$ 对应于 Galois 模 $\sF_{\bar\eta}$，$j_*\sF$ 的特殊纤维则对应 $\sF_{\bar\eta}^I$。

设 $i$ 为闭点 $s$ 的包含，对正合列
\[\xymatrix@C=14pt{
0\ar[r]&j_!\sF\ar[r]&j_*\sF\ar[r]&i_*\sF_{\bar\eta}^I\ar[r]&0\\
0\ar[r]&j_!\sF_{\bar\eta}^I\ar[r]\ar[u]&\sF_{\bar\eta}^I\ar[r]\ar[u]&i_*\sF_{\bar\eta}^I\ar[r]\ar[u]&0
}\]
应用 $D$，得到三角之间的态射
\[\xymatrix@C=14pt{
i_*(\sF_{\bar\eta}^I)^\vee(-1)[-2]\ar[r]\ar@{=}[d]&Dj_*\sF\ar[r]\ar[d]&\eR j_*\sF^\vee\ar[d]\\
i_*(\sF_{\bar\eta}^I)^\vee(-1)[-2]\ar[r]&(\sF_{\bar\eta}^I)^\vee\ar[r]&\eR j_*(\sF_{\bar\eta}^I)^\vee.
}\]
由第一行的长正合列，$i>2$ 时 $\Ext^i=0$；在 $s$ 处取茎，得到
\[
\resizebox{\linewidth}{!}{$\xymatrix@C=12pt{
0\ar[r]&(\underline\h^1(Dj_*\sF))_s\ar[r]&\h^1(I,\sF_{\bar\eta}^\vee)\ar[r]^-\partial\ar[d]&(\sF_{\bar\eta}^I)^\vee(-1)\ar[r]\ar@{=}[d]&(\underline\h^2(Dj_*\sF))_s\ar[r]&0\\
&0\ar[r]&\h^1(I,(\sF_{\bar\eta}^I)^\vee)\ar[r]^-\partial&(\sF_{\bar\eta}^I)^\vee(-1)\ar[r]&0.
}$}
\]
由于 $(\sF_{\bar\eta}^I)^\vee=(\sF_{\bar\eta}^\vee)_I$，置 $M=\sF_{\bar\eta}^\vee$，最后只须验证
\[
\h^1(I,M)\iso\h^1(I,M_I).
\]
若 $p$ 为剩余特征指数，$I$ 是以一个与 $\widehat\dZ_{p'}=\varprojlim_{(m,p)=1}\dZ/m$ 同构的群为商、以一个 $p$-群 $P$ 为核的扩张。由于 $P$ 的阶与 $M$ 的阶互素，$i>0$ 时 $\h^i(P,M)=0$，且 $(M_I)^P=(M^P)_I$。因此可将 $M$ 换成 $M^P$，$I$ 换成 $I/P$。最后，有关于 $M$ 的函子性同构
\[
\h^1(\widehat\dZ_{p'},M)\simeq M_{\widehat\dZ_{p'}},
\]
右边为 $\widehat\dZ_{p'}$ 作用于 $M$ 的\term{余不变量}{coinvariants}{coinvariants}。定理得证。

\begin{theorem_}\label{V:1-4}
对 $\sK\in\ob\eD_c^b(S,\dZ/n)$，有 $\sK\iso DD\sK$。
\end{theorem_}
通过分解归约，可假定 $\sK$ 为集中于次数 $0$ 的可构造层 $\sF$，且 $\sF$ 或有有限支撑（\ref{V:1-2}），或如 \ref{V:1-3} 中那样为 $j_*\sF_1$。后一种情形，由 \ref{V:1-3} 归结为 $U$ 上局部常层 $\sF$ 的局部双对偶。

\section{M.\ Artin 对曲线上 Poincaré 对偶的证明}\label{V:2}
\origpage{164}{158}
\subsection{}\label{V:2-1}
设 $X$ 为代数闭域 $k$ 上的光滑射影曲线。置 $\sK_X=\dZ/n(1)[2]$，并对 $\sK\in\ob\eD_c^b(X,\dZ/n)$ 置 $D\sK=\rHom(\sK,\sK_X)$。对 $\dZ/n$-模 $M$，也记 $DM=\hom(M,\dZ/n)$；模复形的情形亦同。迹态射 $\h^0(X,\sK_X)\to\dZ/n$，或等价地 $\eR\Gamma(X,\sK_X)\to\dZ/n$，定义配对
$\eR\Gamma(X,\sK)\lotimes\eR\Gamma(X,D\sK)\to\dZ/n$。$X$ 的开集 $j:U\hookrightarrow X$ 上，普通上同调与固有支撑上同调之间的 Poincaré 对偶，正是说当 $\sK=j_!\dZ/n$ 时，此配对将每个因子识别为另一个因子的对偶。

下面介绍 M.\ Artin 告诉我的证明：对任意 $\sK\in\ob\eD_c^b(X,\dZ/n)$，这个配对总是\term{完美配对}{perfect pairing}{accouplement parfait}，即定义同构
\begin{equation*}\tag{2.1.1}\label{V:eq:2-1-1}
\eR\Gamma(X,\sK)\iso D\eR\Gamma(X,D\sK).
\end{equation*}
对每个可构造层 $\sF$，置
\[
'\h^i(X,\sF)=\hom(\h^{-i}(X,D\sF),\dZ/n).
\]
\eqref{V:eq:2-1-1} 为同构，等价于下述定理。

\begin{theorem_}\label{V:2-2}
对任意可构造 $\dZ/n$-模层 $\sF$，有
\begin{equation*}\tag{2.2.1}\label{V:eq:2-2-1}
\h^i(X,\sF)\iso'\h^i(X,\sF).
\end{equation*}
\end{theorem_}

\begin{lemma_}\label{V:2-3}
设 $f:X\to Y$ 为 $k$ 上光滑曲线之间泛平展的态射。则 $\sK_X=f^*\sK_Y=\eR f^!\sK_Y$，而伴随箭头 $\eR f_!\eR f^!\sK_Y\to\sK_Y$ 是 $\tr_f$。
\end{lemma_}
可归结为验证 $f$ 有限的情形：
\[
f_*\eR f^!\sK_Y=\rHom(f_*\dZ/n,\sK_Y)\simeq f_*\dZ/n(1)[2],
\]
伴随箭头为 $\tr_f$。$i>0$ 时 $\Ext^i$ 的消失由 \ref{V:1-3} 得到，断言随之成立。\translatornote{原书这一证明句将右端简写为 $f_*\dZ/n$，漏去了两边共有的扭曲和平移 $(1)[2]$。保留 $\sK_Y=\dZ/n(1)[2]$ 时，应为这里写出的同构；去掉 $(1)[2]$ 后，所需断言才是 $\Hom(f_*\dZ/n,\dZ/n)\simeq f_*\dZ/n$。后者的配对在平展处是各分支上乘积的和，因而与迹相符，并由 1.3 延拓。}

更明确地，$\eR f^!$ 是函子 $f^!$ 的导出函子，其中 $f^!\sF(U)=\hom(f_!\dZ_U,\sF)$；而 \ref{V:1-3} 蕴含 $i>0$ 时 $\eR^if^!\dZ/n=0$。

\begin{corollary_}\label{V:2-4}
设 $f:X'\to X$ 为 $k$ 上光滑射影曲线之间的泛平展态射。则 $'\h^i(X,f_*\sF)='\h^i(X',\sF)$；该同构具函子性，且下图交换：
\[\xymatrix{\h^i(X,f_*\sF)\ar[r]\ar@{=}[d]&'\h^i(X,f_*\sF)\ar@{=}[d]\\
\h^i(X',\sF)\ar[r]&'\h^i(X',\sF).}\]
\end{corollary_}
这个同构由 \ref{V:1-2} 给出：$Df_*\sF=f_*D\sF$。

\subsection{}\label{V:2-5}
证明 \ref{V:2-2}。若 $\sF$ 是一个点上的 $\dZ/n$，再以零延拓，则 $\h^\bullet(D\sF)=\ext^\bullet(\sF,\sK_X)$ 为集中于次数 $0$ 的 $\dZ/n$，而配对 $\h^\bullet(\sF)\otimes\h^\bullet(D\sF)\to\dZ/n$ 将 $1\otimes1$ 送到 $1$（\hyperref[IV]{闭链}，\ref{IV:2-1-5}），故此对偶完美。$\sF$ 有有限支撑时同理。

对任意可构造层 $\sF$，记 $\sA$ 为其有限支撑截面组成的子层，$\sG=\sF/\sA$。正合列
\begin{equation*}\tag{2.5.1}\label{V:eq:2-5-1}
\xymatrix{0\ar[r]&\sA\ar[r]&\sF\ar[r]&\sG\ar[r]&0}
\end{equation*}
给出图
\begin{equation*}\tag{2.5.2}\label{V:eq:2-5-2}
\resizebox{\linewidth}{!}{$\xymatrix@C=12pt{
&&0\ar[r]&\h^0(X,\sA)\ar[r]\ar[d]^-\wr&\h^0(X,\sF)\ar[r]\ar[d]&\h^0(X,\sG)\ar[r]\ar[d]&0\\
0\ar[r]&'\h^{-1}(X,\sF)\ar[r]&'\h^{-1}(X,\sG)\ar[r]&'\h^0(X,\sA)\ar[r]&'\h^0(X,\sF)\ar[r]&'\h^0(X,\sG)\ar[r]&0
}$}
\end{equation*}
以及彼此相容的同构 $\h^i(X,\sF)\iso\h^i(X,\sG)$、$'\h^i(X,\sF)\iso'\h^i(X,\sG)$（$i\ne0,-1$）。取开集 $j:U\to X$ 使 $\sF$ 在其上局部常，便有正合列
\begin{equation*}\tag{2.5.3}\label{V:eq:2-5-3}
\xymatrix{0\ar[r]&\sG\ar[r]&j_*j^*\sF\ar[r]&\sB\ar[r]&0,}
\end{equation*}
其中 $\sB$ 有有限支撑。由 \ref{V:1-3}，
$'\h^i(X,j_*j^*\sF)=D\h^{2-i}(X,j_*(j^*\sF)^\vee(1))$，故 $i<0$ 时此群为零。结合 \eqref{V:eq:2-5-2} 和 \eqref{V:eq:2-5-3} 的长正合列，得到 $i<0$ 时 $'\h^i$ 总为零。

若 $\sF=\dZ/n$，定理在 $i=0,2$ 时成立；这就是说，$X$ 连通时，迹态射 $\h^2(X,\dmu_n)\iso\dZ/n$ 是同构。由 \ref{V:2-4}，定理对 $i=0$ 及 $\sF=f_*\dZ/n$ 仍成立。

再令 $\sF$ 任意，$\sG$ 如上。存在 \ref{V:2-4} 中的 $f:X'\to X$，使 $\sG$ 嵌入 $f_*\dZ/n$：
\[\xymatrix{0\ar[r]&\sG\ar[r]&f_*\dZ/n\ar[r]&\sQ\ar[r]&0.}\]
将 $X'$ 换成其上的 $X''/X'$ 后，可以假定对 $i>0$，映射
$\h^i(X,\sG)\to\h^i(X,f_*\dZ/n)=\h^i(X',\dZ/n)$ 为零。于是考虑
\[
\resizebox{\linewidth}{!}{$\xymatrix@C=10pt{
0\ar[r]&\h^0(X,\sG)\ar[r]\ar[d]&\h^0(X',\dZ/n)\ar[r]\ar[d]^-\wr&\h^0(X,\sQ)\ar[r]\ar[d]&\h^1(X,\sG)\ar[r]^-0\ar[d]&\h^1(X',\dZ/n)\ar[r]\ar[d]&\h^1(X,\sQ)\ar[r]\ar[d]&\cdots\\
0\ar[r]&'\h^0(X,\sG)\ar[r]&'\h^0(X',\dZ/n)\ar[r]&'\h^0(X,\sQ)\ar[r]&'\h^1(X,\sG)\ar[r]&'\h^1(X',\dZ/n)\ar[r]&'\h^1(X,\sQ)\ar[r]&\cdots
}$}
\]
$\h^0(X')$ 处的同构先给出 $\h^0(X,\sG)\hookrightarrow'\h^0(X,\sG)$；由 \eqref{V:eq:2-5-2}，对 $\sF$ 也得到同样的单射性。将这一结论用于 $\sQ$，即可得 $\h^0(X,\sG)\iso'\h^0(X,\sG)$；再由 \eqref{V:eq:2-5-2}，对 $\sF$ 亦然。再次将它用于 $\sQ$，得到 $\h^1(X,\sG)\hookrightarrow'\h^1(X,\sG)$。把此单射性用于 $\h^1(X',\dZ/n)$，由于该群与 $'\h^1(X',\dZ/n)=D(\h^1(X',\dZ/n)(1))$ 阶数相同，便得到同构，从而 $\h^1(X,\sG)\iso'\h^1(X,\sG)$；$\sF$ 的情形也一样。对 $i=2$，已经知道 $\h^2(X',\dZ/n)\iso'\h^2(X',\dZ/n)$，照同样方法即可。再往后，各项都已为零，证明至此结束。

\section{曲线上 Poincaré 对偶的相容性}\label{V:3}
\origpage{167}{161}
本节证明（\hyperref[I]{Arcata}，\ref{I:eq:6-2-3-3}）中的相容性。

\subsection{}\label{V:3-1}
沿用（\hyperref[I]{Arcata}，\ref{I:6-2-3}）的记号：$\bar X$ 为代数闭域 $k$ 上的光滑连通射影曲线，$D$ 为 $\bar X$ 上的约化除子，$0$ 为 $X=\bar X\setminus D$ 的一个点，
\[\xymatrix{X\ar@{^{(}->}[r]^-j&\bar X&\ar[l]_-iD,}\]
且 $\pic_D(\bar X)=\h^1(\bar X,{}_D\dG_m)$，其中 ${}_D\dG_m$ 为 $\dG_m$ 中模 $D$ 同余于 $1$ 的截面组成的层。回顾
$\h_c^1(X,\dmu_n)=\pic_D^0(\bar X)_n$，且 $x\mapsto\sO(x)$ 给出典范映射 $f:X\to\pic_D(\bar X)$。置 $f_0(x)=f(x)-f(0):X\to\pic_D^0(\bar X)$。

仍用 $j$ 表示 $X\times X$ 到 $X\times\bar X$ 的包含。$X\times X$ 的对角线 $\Delta$ 闭于 $X\times\bar X$，因而定义类
$\h_\Delta^2(X\times X,\dmu_n)=\h_\Delta^2(X\times\bar X,j_!\dmu_n)$。记它在 $\h^2(X\times\bar X,j_!\dmu_n)$ 中的像为 $c$。

由 Künneth 公式，
\[
\begin{aligned}
\h^\bullet(X\times\bar X,j_!\dmu_n)
&=\h^\bullet(X,\dZ/n)\otimes\h^\bullet(\bar X,j_!\dmu_n)\\
&=\h^\bullet(X,\dZ/n)\otimes\h_c^\bullet(X,\dmu_n).
\end{aligned}
\]
我们要计算 $c$ 的 $(1,1)$-分量
\[
\begin{aligned}
c^{1,1}&\in\h^1(X,\dZ/n)\otimes\h_c^1(X,\dmu_n)\\
&=\h^1(X,\h_c^1(X,\dmu_n))=\h^1(X,\pic_D^0(\bar X)_n).
\end{aligned}
\]
正合列
\[\xymatrix@C=15pt{0\ar[r]&\pic_D^0(\bar X)_n\ar[r]&\pic_D^0(\bar X)\ar[r]^-n&\pic_D^0(\bar X)\ar[r]&0}\]
使 $\pic_D^0(\bar X)$ 成为 $\pic_D^0(\bar X)$ 上的 $\pic_D^0(\bar X)_n$-扭子。将它沿 $f_0$ 拉回，记所得扭子在
$\h^1(X,\h_c^1(X,\dmu_n))=\h^1(X,\pic_D^0(\bar X)_n)$ 中的类为 $u$。

\begin{proposition_}\label{V:3-2}
在上述记号下，$c^{1,1}=-u$。
\end{proposition_}
设 $e$ 为 $\Delta$ 和 $X\times\{0\}$ 的上同调类在 $\h^2(X\times\bar X,j_!\dmu_n)$ 中的像之差。$X\times\{0\}$ 的类为 Künneth 型 $(0,2)$，故 $c^{1,1}=e^{1,1}$。对 $\pr_1:X\times\bar X\to X$ 及层 $j_!\dmu_n$ 取 Leray 谱序列：
\[
E_2^{pq}=\h^p(X,\eR^q\pr_{1*}j_!\dmu_n)
=\h^p(X,\h_c^q(X,\dmu_n))
\Rightarrow\h^{p+q}(X\times\bar X,j_!\dmu_n).
\]
它退化，因为它是 $\h^\bullet(X)$ 与 $\bar X\to\spec(k)$、层 $j_!\dmu_n$ 的平凡 Leray 谱序列的张量积。除子 $\Delta-X\times\{0\}$ 在每条纤维上的次数均为 $0$，故 $e$ 在 $E^{02}$ 中的像为零。因此可以考虑它在 $E^{1,1}$ 中的像 $\bar e$；只须证明 $\bar e=-u$。

设 $\sG$ 为 $X\times\bar X$ 上 $\dG_m$ 的子层，由限制到子概形 $X\times D$ 后等于 $1$ 的局部截面组成。仍有正合列
\begin{equation*}\tag{3.2.1}\label{V:eq:3-2-1}
\xymatrix{0\ar[r]&j_!\dmu_n\ar[r]&\sG\ar[r]^-n&\sG\ar[r]&0,}
\end{equation*}
而 $e$ 是 $\sO(\Delta-X\times\{0\})$ 连同其由 $1$ 给出的 $X\times D$ 上平凡化的类
$e_1\in\h^1(X\times\bar X,\sG)$ 经 $\partial$ 的像。

由 \eqref{V:eq:3-2-1} 定义的特异三角，经 $\eR\pr_{1*}$ 直接像后仍为特异三角：
\begin{equation*}\tag{3.2.2}\label{V:eq:3-2-2}
\eR\pr_{1*}j_!\dmu_n\longrightarrow\eR\pr_{1*}\sG
\xrightarrow{n}\eR\pr_{1*}\sG
\xrightarrow{\partial}(\eR\pr_{1*}j_!\dmu_n)[1].
\end{equation*}
下图交换：
\[\xymatrix@C=16pt{
\h^1(X\times\bar X,\sG)\ar[r]^-\partial\ar@{=}[d]^-{(1)}&\h^2(X\times\bar X,j_!\dmu_n)\ar@{=}[d]\\
\h^1(X,\eR\pr_{1*}\sG)\ar[r]^-\partial&\h^2(X,\eR\pr_{1*}j_!\dmu_n).
}\]
仍记由 (1) 对应于 $e_1$ 的类为 $e_1$。问题于是变为：求它经
$\partial_{\eqref{V:eq:3-2-2}}:\h^1(X,\eR\pr_{1*}\sG)\to\h^2(X,\eR\pr_{1*}j_!\dmu_n)$ 后，在 $E^{11}=\h^1(X,\h_c^1(X,\dmu_n))$ 中的像（其在 $E^{02}$ 中的像为零）。

层 $\eR^1\pr_{1*}\sG$ 由 $X$ 上的群概形表示，该群概形是 $\spec(k)$ 上代数群 $\pic_D(\bar X)$ 的基变换。$e_1$ 在
$\h^0(X,\eR^1\pr_{1*}\sG)=\hom(X,\pic_D(\bar X))$ 中的像为 $f_0$。

用层复形的短正合列
\[\xymatrix{0\ar[r]&\sA\ar[r]&\sB\ar[r]&\sC\ar[r]&0}\]
表示三角 \eqref{V:eq:3-2-2}。在截断复形
\[
\tau_{\leqslant1}\sB:\quad\cdots\longrightarrow\sB^0\longrightarrow\ker(d)\longrightarrow0\longrightarrow\cdots
\]
中，将一次项 $\ker(d)$ 换成 $X$ 上 $\pic_D^0(\bar X)$ 在
$\ker(d)\to\underline\h^1(\sB)=\pic_D(\bar X)$ 下的逆像，所得子复形记为 $\sB'$。置 $\sA'=\sA\cap\sB'$，并记 $\sB'$ 的像为 $\sC'$。则 $\sA'=\tau_{\leqslant1}\sA$，而 $\sC'$ 可由 $\sC$ 作与 $\sB'$ 之于 $\sB$ 相同的构造得到。于是得到短正合列的交换图
\begin{equation*}\tag{3.2.3}\label{V:eq:3-2-3}
\xymatrix@C=15pt{
0\ar[r]&\sA\ar[r]&\sB\ar[r]&\sC\ar[r]&0\\
0\ar[r]&\tau_{\leqslant1}\sA\ar[r]\ar[u]\ar[d]&\sB'\ar[r]\ar[u]\ar[d]&\sC'\ar[r]\ar[u]\ar[d]&0\\
0\ar[r]&\h_c^1(X,\dmu_n)[-1]\ar[r]&\pic_D^0(\bar X)[-1]\ar[r]^-n&\pic_D^0(\bar X)[-1]\ar[r]&0,
}
\end{equation*}
其中最后一行应视为 $X$ 上集中于次数 $1$ 的层复形的正合列。\translatornote{原书此图将集中于次数 $1$ 的层写成 $F[1]$，应为 $F[-1]$，因为上同调平移满足 $H^i(F[r])=H^{i+r}(F)$。其后紧接着使用的 $\hh^2(X,F[-1])=\h^1(X,F)$ 也确认了这一点。原书还把此图重复标为 (3.2.2)；译本顺次编号为 (3.2.3)，并将前一三角中误写的 $\pr_{2*}$ 统一为 $\pr_{1*}$。}
\begin{enumerate}[a)]
\item $e\in\hh^2(X,\sA)$ 来自某个 $\widetilde e\in\hh^2(X,\tau_{\leqslant1}\sA)$，因为它在 $E^{02}$ 中的像为零。所求的是 $\widetilde e$ 在
$\hh^2(X,\h_c^1(X,\dmu_n)[-1])=\h^1(X,\h_c^1(X,\dmu_n))$ 中的像 $\bar e$。
\item $e_1\in\hh^1(X,\sC)$ 来自某个 $\widetilde e_1\in\hh^1(X,\sC')$，因为它在 $\h^0(X,\pic_D(\bar X))$ 中的像 $f_0$ 属于 $\h^0(X,\pic_D^0(\bar X))$。可以取 $\widetilde e=\partial\widetilde e_1$。
\item 因而 $\bar e=\partial f_0$，这里 $\partial$ 来自上述复形正合列，并使用通常的同构 $\hh^2(X,\sF[-1])=\h^1(X,\sF)$。这些同构与连接同态 $\partial$ 反交换，因此若将 $\partial$ 改指正合列
\[\xymatrix@C=15pt{0\ar[r]&\h_c^1(X,\dmu_n)\ar[r]&\pic_D^0(\bar X)\ar[r]^-n&\pic_D^0(\bar X)\ar[r]&0}\]
的连接同态，就有 $\bar e=-\partial f_0$。由（\hyperref[IV]{闭链}，\ref{IV:1-1-1}）即得结论。
\end{enumerate}

\subsection{}\label{V:3-3}
对任意同态 $\varphi:\h_c^1(X,\dmu_n)\to\dZ/n$，记 $u$ 经
$\varphi:\h^1(X,\h_c^1(X,\dmu_n))\to\h^1(X,\dZ/n)$ 的像为 $\varphi(u)$。若 $x\in\h_c^1(X,\dmu_n)$，则杯积 $\varphi(u)\smallsmile x$ 属于 $\h_c^2(X,\dmu_n)$。我们要证明：

\begin{proposition_}\label{V:3-4}
$\tr(\varphi(u)\smallsmile x)=\varphi(x)$。
\end{proposition_}
该恒等式可由
\begin{equation*}\tag{3.4.1}\label{V:eq:3-4-1}
\tr(u\smallsmile x)=x
\end{equation*}
应用 $\varphi$ 得到；这里的 $\tr$ 是映射
\[
\tr:\h_c^2(X,\h_c^1(X,\dmu_n)\otimes\dmu_n)\longrightarrow\h_c^1(X,\dmu_n).
\]
因此 $x\mapsto\tr(u\smallsmile x)$ 是以下复合：
\begin{equation*}\tag{3.4.2}\label{V:eq:3-4-2}
\begin{aligned}
\h_c^1(X,\dmu_n)&\xrightarrow{x\mapsto u\otimes x}
\h^1(X,\h_c^1(X,\dmu_n))\otimes\h_c^1(X,\dmu_n)\\
&\xleftarrow{\sim}\h^1(X,\dZ/n)\otimes\h_c^1(X,\dmu_n)\otimes\h_c^1(X,\dmu_n)\\
&\xrightarrow{(1)}\h_c^2(X,\dmu_n)\otimes\h_c^1(X,\dmu_n)
\xrightarrow{\tr\otimes\mathrm{id}}\h_c^1(X,\dmu_n),
\end{aligned}
\end{equation*}
其中 (1) 由 $a\otimes b\otimes c\mapsto(a\smallsmile c)\otimes b$ 给出。

将这个态射用相对于第二投影的迹态射
\[
\tr:\h_c^3(X\times X,\dmu_n^{\otimes2})\longrightarrow\h_c^1(X,\dmu_n)
\]
表示。对 $x\in\h_c^3$，计算 $\tr(x)$ 时，只保留其 Künneth 分解中型为 $(2,1)$ 的分量，然后对第一个因子应用 $\tr:\h_c^2(X,\dmu_n)\to\dZ/n$。

记 $u$ 经下列映射的像为 $u'$：
\[
\begin{aligned}
\h^1(X,\h_c^1(X,\dmu_n))
&\xleftarrow{\sim}\h^1(X,\dZ/n)\otimes\h_c^1(X,\dmu_n)\\
&=\h^1(X,\dZ/n)\otimes\h^1(\bar X,j_!\dmu_n)\\
&\longrightarrow\h^2(X\times\bar X,\dZ/n\boxtimes j_!\dmu_n).
\end{aligned}
\]
由 Künneth 同构与杯积的关系，
\[
\tr(u\smallsmile x)=-\tr(u'\smallsmile\pr_1^*x).
\]
第二个杯积是（\hyperref[IV]{闭链}，\ref{IV:1-2-4}）中的杯积：
\[
\begin{aligned}
&\h^2(X\times\bar X,\dZ/n\boxtimes j_!\dmu_n)
\otimes\h^1(\bar X\times\bar X,j_!\dmu_n\boxtimes\dZ/n)\\
&\hspace{3em}\longrightarrow\h^3(\bar X\times\bar X,j_!\dmu_n\boxtimes j_!\dmu_n).
\end{aligned}
\]
负号来自 (1) 中两个次数为 $1$ 的符号交换。$c\smallsmile\pr_1^*x$ 与 $c^{1,1}\smallsmile\pr_1^*x$ 有相同的 Künneth $(2,1)$-分量，又 $u'=-c^{1,1}$，故
\[
\tr(u\smallsmile x)=\tr(c\smallsmile\pr_1^*x).
\]
\translatornote{原书上一句印作 $u'=c^{1,1}$，应补负号。命题 3.2 已证明 $c^{1,1}=-u$，而 $u'$ 就是 $u$ 经 Künneth 同构的像，所以 $u'=-c^{1,1}$。这个负号恰与交换两个一次因子所产生的负号抵消，才得到此处正号的等式。}

右侧杯积这次可解释为（\hyperref[IV]{闭链}，\ref{IV:1-2-5}）中的杯积：将
$\cl(\Delta)\in\h_\Delta^2(X\times X,\dmu_n)$ 与
$\pr_1^*x\in\h_c^1(\Delta,\dmu_n)$ 相乘。在 $\Delta$ 上，$\pr_1=\pr_2$，所以
\[
\begin{aligned}
\tr(c\smallsmile\pr_1^*x)
&=\tr(\cl(\Delta)\smallsmile\pr_1^*x)\\
&=\tr(\cl(\Delta)\smallsmile\pr_2^*x)
=\tr(\cl(\Delta))\smallsmile x=x.
\end{aligned}
\]

\section{SGA 4 XVIII 3.1.10.3 的相容性}\label{V:4}
\origpage{172}{166}
在 \cite[XVIII]{sga4} 中，我一时泄气，没有核实标题所指的相容性。虽然现在看来它不过是虚惊一场，我仍觉得有责任补上这个缺口。对分离有限型态射 $f:X\to S$，须验证伴随态射 $\eR f_!\eR f^!\sL\to\sL$ 与任意平展局部化 $k:V\to S$ 相容。如下图，以 $k,f$ 分别表示若干箭头，所需验证的是右图的交换性：
\[
\xymatrix{X_V\ar[r]^-k\ar[d]^-f&X\ar[d]^-f\\V\ar[r]^-k&S}
\qquad
\xymatrix@C=16pt{k^*\eR f_!\eR f^!\sL\ar[r]^-{k^*(\mathrm{adj})}\ar[d]^-\sim&k^*\sL\ar@{=}[d]\\
\eR f_!\eR f^!k^*\sL\ar[r]^-{\mathrm{adj}}&k^*\sL.}
\]
在复形层次（同上所引文献），所需交换性写为
\[\xymatrix{k^*f_!^\bullet f^!_\bullet\sL\ar[r]\ar[d]^-{(1)}&k^*\sL\ar@{=}[d]\\
f_!^\bullet f^!_\bullet k^*\sL\ar[r]&k^*\sL.}\]
逐项验证，即归结为 $\sL$ 是层时，对每对伴随函子 $(f_i^!,f_!^i)$ 的相应相容性；简记它们为 $f^!,f_!$。箭头 (1) 由同构 $(2):k^*f_!=f_!k^*$ 及其诱导的态射 $(3):k^*f^!\to f^!k^*$ 定义。

在所引文献中，先利用 $k^*=k^!$ 与 $k_!$ 的伴随关系，从 (2) 定义态射 $(4):k_!f_!\to f_!k_!$：
\[
k_!f_!\longrightarrow k_!f_!k^*k_!
\longrightarrow k_!k^*f_!k_!\longrightarrow f_!k_!,
\]
再由 (4) 经伴随得到 (3)。等价地，也可以直接利用 $f_!$ 与 $f^!$ 的伴随关系从 (2) 得到 (3)：
\[
k^*f^!\longrightarrow f^!f_!k^*f^!
\longrightarrow f^!k^*f_!f^!\longrightarrow f^!k^*.
\]

将 $k^*$ 简写为 $k$，并略去 $\sL$，所求就是下图外边界的交换性：
\[\xymatrix@C=12pt@R=18pt{
kf_!f^!\ar[r]\ar@{=}[d]&f_!kf^!\ar[r]&f_!f^!f_!kf^!\ar[r]&f_!f^!kf_!f^!\ar@{=}[d]\\
kf_!f^!\ar[d]&&&\ar[lll]f_!f^!kf_!f^!\ar[d]\\
k&&&\ar[lll]f_!f^!k.
}\]
第一行给出同态 $kf_!f^!\to f_!f^!kf_!f^!$；它的定义利用了 $kf_!f^!$ 可写成 $f_!X$ 的形式（$X=kf^!$）。第二行的伴随态射是它的回缩；下方方块的交换性则正是该伴随态射的函子性。
